%% paper69_hodge_vermutung_de.tex
%% Paper 69 (DE): Die Hodge-Vermutung — Ein Millennium-Problem an der Schnittstelle
%%                von Topologie, Algebra und komplexer Geometrie
%% Autor: Michael Fuhrmann
%% Projekt: Specialist Mathematics, Build 147
%% Datum: 2026-03-12

\documentclass[12pt,a4paper]{amsart}
\usepackage{amsmath,amssymb,amsthm}
\usepackage{mathtools}
\usepackage[utf8]{inputenc}
\usepackage[T1]{fontenc}
\usepackage[ngerman]{babel}
\usepackage{hyperref}
\usepackage{enumitem,booktabs,array}
\usepackage{geometry}
\geometry{margin=2.5cm}

\newtheorem{theorem}{Satz}[section]
\newtheorem{lemma}[theorem]{Lemma}
\newtheorem{korollar}[theorem]{Korollar}
\newtheorem{proposition}[theorem]{Proposition}
\newtheorem{vermutung}[theorem]{Vermutung}
\theoremstyle{definition}
\newtheorem{definition}[theorem]{Definition}
\newtheorem{beispiel}[theorem]{Beispiel}
\theoremstyle{remark}
\newtheorem{bemerkung}[theorem]{Bemerkung}

\title{Die Hodge-Vermutung:\\
Ein Millennium-Problem an der Schnittstelle von Topologie, Algebra und komplexer Geometrie}

\author{Michael Fuhrmann}
\address{Specialist Mathematics Project, Build 147}
\email{specialist-maths@localhost}
\date{\today}

\begin{abstract}
Die Hodge-Vermutung ist eines der sieben Millennium-Preisprobleme des Clay Mathematics
Institute (2000), für deren Lösung jeweils eine Million US-Dollar ausgesetzt ist.
Sie behauptet, dass auf einer glatten projektiven komplexen algebraischen Varietät
jede rationale Kohomologieklasse vom Hodge-Typ $(p,p)$ eine rationale Linearkombination
von Klassen algebraischer Teilvarietäten (algebraischer Zyklen) ist. Diese Arbeit
gibt eine in sich abgeschlossene Einführung in die mathematische Landschaft rund um
die Vermutung: Kähler-Mannigfaltigkeiten und die Kähler-Form $\omega$, das
$\partial\bar\partial$-Lemma, der Hodge-Zerlegungssatz, de-Rham- und Dolbeault-Kohomologie,
algebraische Zyklen und ihre Fundamentalklassen sowie die formale Aussage der Vermutung.
Wir diskutieren den Lefschetz-$(1,1)$-Satz — die einzige nicht-triviale Dimension,
in der die Hodge-Vermutung vollständig bewiesen ist — sowie Gegenbeispiele, die zeigen,
dass weder die ganzzahlige Version noch das Analogon für Nicht-Kähler-Mannigfaltigkeiten
gilt. Die Rolle der Grothendieck'schen Theorie der Motive und der Standard-Vermutungen
als konjekctureller Rahmen wird skizziert. Durchgehend werden bewiesene Sätze und
offene Vermutungen sorgfältig unterschieden, da die Hodge-Vermutung selbst vollständig
offen ist.
\end{abstract}

\maketitle

\tableofcontents

%% ============================================================
\section{Einleitung}
\label{sec:einleitung}
%% ============================================================

Im Jahr 1950, beim Internationalen Mathematiker-Kongress in Cambridge (Massachusetts),
stellte William Vallance Douglas Hodge ein weitreichendes Programm vor, das die
Topologie komplexer algebraischer Varietäten mit ihrer algebraischen Geometrie
verbindet. Im Mittelpunkt stand, was schließlich als \emph{Hodge-Vermutung} bekannt
werden sollte:

\begin{quotation}
\itshape
Auf einer glatten projektiven komplexen algebraischen Varietät $X$ ist jede rationale
Kohomologieklasse vom Typ $(p,p)$ eine rationale Linearkombination von Klassen
algebraischer Zyklen.
\end{quotation}

Ein halbes Jahrhundert später nahm das Clay Mathematics Institute sie in die Liste
der sieben \emph{Millennium-Preisprobleme} auf \cite{clay2000}. Trotz enormer
Fortschritte in algebraischer Geometrie, Hodge-Theorie und der Theorie der Motive
in den seitdem vergangenen Jahrzehnten ist die Vermutung in allen Fällen jenseits
der komplexen Dimension~1 ungelöst.

\subsection*{Warum die Hodge-Vermutung wichtig ist}

Die Vermutung steht an einer bemerkenswerten Schnittstelle mathematischer Disziplinen:

\begin{enumerate}[label=(\roman*)]
\item \textbf{Topologie}: Kohomologieklassen sind topologische Invarianten, ohne
  Bezug auf algebraische oder komplexe Strukturen berechnet.
\item \textbf{Komplexe Geometrie}: Die Hodge-Zerlegung $H^n(X,\mathbb{C}) =
  \bigoplus_{p+q=n} H^{p,q}(X)$ hängt von der komplexen Struktur ab und ist das
  zentrale Objekt der Hodge-Theorie.
\item \textbf{Algebraische Geometrie}: Algebraische Zyklen — durch Polynomgleichungen
  definierte Teilvarietäten — erzeugen Kohomologieklassen, die konstruktionsbedingt
  vom Typ $(p,p)$ sind.
\end{enumerate}

Die Vermutung behauptet, dass diese drei Perspektiven in einem präzisen Sinn
\emph{kompatibel} sind: Jede kohomologisch nachweisbare $(p,p)$-Klasse entsteht
geometrisch aus einer algebraischen Teilvarietät.

\subsection*{Gliederung}

Abschnitt~\ref{sec:kaehler} führt Kähler-Mannigfaltigkeiten ein. Abschnitt~\ref{sec:hodge-theorie}
entwickelt den Hodge-Satz via Laplace–Beltrami-Operator. Abschnitt~\ref{sec:kohomologie}
vergleicht de-Rham- und Dolbeault-Kohomologie. Abschnitt~\ref{sec:zyklen} behandelt
algebraische Zyklen. Abschnitt~\ref{sec:vermutung} formuliert die Hodge-Vermutung
präzise. Abschnitt~\ref{sec:bekannt} gibt eine Übersicht bekannter Fälle.
Abschnitt~\ref{sec:lefschetz} behandelt den Lefschetz-$(1,1)$-Satz ausführlich.
Abschnitt~\ref{sec:gegenbeispiele} diskutiert Gegenbeispiele. Abschnitt~\ref{sec:motive}
skizziert den motivischen Rahmen. Abschnitt~\ref{sec:literatur} enthält Literaturangaben.

%% ============================================================
\section{Kähler-Mannigfaltigkeiten}
\label{sec:kaehler}
%% ============================================================

\subsection{Hermitesche Metriken und die Kähler-Form}

Sei $X$ eine kompakte komplexe Mannigfaltigkeit der komplexen Dimension $n$. Das
Tangentialbündel $TX$ trägt eine natürliche fast-komplexe Struktur $J: TX \to TX$
mit $J^2 = -\mathrm{id}$.

\begin{definition}
Eine \emph{hermitesche Metrik} auf $X$ ist eine $J$-kompatible Riemannsche Metrik
$g$, d.h. $g(Ju, Jv) = g(u,v)$ für alle Tangentialvektoren $u,v$. Die zugehörige
\emph{fundamentale (Kähler-)Form} ist die reelle $(1,1)$-Form
\[
  \omega(u,v) := g(Ju, v).
\]
\end{definition}

In lokalen holomorphen Koordinaten $(z^1, \ldots, z^n)$ schreibt man
$\omega = \frac{i}{2}\sum_{j,k} h_{j\bar k}\, dz^j \wedge d\bar z^k$,
wobei $(h_{j\bar k})$ eine positiv-definite hermitesche Matrix ist.

\begin{definition}
Eine hermitesche Mannigfaltigkeit $(X, g, \omega)$ heißt \emph{Kähler-Mannigfaltigkeit},
wenn $d\omega = 0$, d.h. die Fundamentalform ist geschlossen.
\end{definition}

Diese eine Bedingung — die Geschlossenheit von $\omega$ — hat weitreichende Folgen.
Sie gewährleistet die Kompatibilität von komplexer, Riemannscher und symplektischer
Struktur.

\begin{beispiel}
\begin{enumerate}[label=(\alph*)]
\item Der projektive Raum $\mathbb{CP}^n$ mit der Fubini–Study-Metrik ist Kähler.
\item Jede glatte projektive Varietät $X \subset \mathbb{CP}^N$ erbt eine Kähler-Struktur
  von der umgebenden Fubini-Study-Metrik.
\item Komplexe Tori $\mathbb{C}^n / \Lambda$ sind Kähler.
\item $K3$-Flächen sind Kähler (Siu, 1983, bewiesen).
\end{enumerate}
\end{beispiel}

\subsection{Das \texorpdfstring{$\partial\bar\partial$}{deldel-quer}-Lemma}

Auf einer kompakten Kähler-Mannigfaltigkeit erfüllen die Operatoren $\partial$,
$\bar\partial$ und ihre Adjungierten $\partial^*$, $\bar\partial^*$ die entscheidende
Identität
\[
  \Delta_d = 2\Delta_\partial = 2\Delta_{\bar\partial},
\]
wobei $\Delta_d = dd^* + d^*d$ der de-Rham-Laplacian und $\Delta_\partial =
\partial\partial^* + \partial^*\partial$ der Dolbeault-Laplacian ist. Dies führt zu:

\begin{theorem}[$\partial\bar\partial$-Lemma, bewiesen]
\label{thm:ddbar}
Sei $X$ eine kompakte Kähler-Mannigfaltigkeit. Wenn eine Form $\alpha$ $d$-geschlossen,
$\partial$-exakt und $\bar\partial$-exakt ist, dann gilt $\alpha = \partial\bar\partial\beta$
für eine Form $\beta$.
\end{theorem}

\begin{proof}[Beweisidee]
Dies folgt aus den Kähler-Identitäten $[\Lambda, \partial] = -i\bar\partial^*$ und
$[\Lambda, \bar\partial] = i\partial^*$, wobei $\Lambda$ der Adjungierte des
Lefschetz-Operators $L(\alpha) = \omega \wedge \alpha$ ist. Diese Identitäten implizieren
$\Delta_d = 2\Delta_\partial = 2\Delta_{\bar\partial}$, woraus das Lemma durch ein
Hodge-Theorie-Argument folgt. Siehe \cite[Kap.~6]{griffiths_harris1978}.
\end{proof}

Das $\partial\bar\partial$-Lemma ist das entscheidende analytische Werkzeug, das die
Kähler-Geometrie so viel reicher macht als die allgemeine komplexe Geometrie.

\subsection{Der harte Lefschetz-Satz}

\begin{theorem}[Harter Lefschetz-Satz, bewiesen]
\label{thm:hard-lefschetz}
Sei $X$ eine kompakte Kähler-Mannigfaltigkeit der komplexen Dimension $n$. Die Abbildung
\[
  L^{n-k} : H^k(X,\mathbb{R}) \longrightarrow H^{2n-k}(X,\mathbb{R}),
  \quad [\alpha] \mapsto [\omega^{n-k} \wedge \alpha]
\]
ist ein Isomorphismus für alle $0 \le k \le n$.
\end{theorem}

Dieser Satz, von Hodge mit harmonischen Formen bewiesen, hat tiefgreifende Konsequenzen
für die Topologie projektiver Varietäten.

%% ============================================================
\section{Hodge-Theorie}
\label{sec:hodge-theorie}
%% ============================================================

\subsection{Der Laplace–Beltrami-Operator}

Sei $(X,g)$ eine kompakte orientierte Riemannsche Mannigfaltigkeit der reellen Dimension $m$.
Bezeichne mit $\Omega^k(X)$ den Raum der glatten $k$-Formen. Das $L^2$-Skalarprodukt ist
\[
  \langle \alpha, \beta \rangle = \int_X \alpha \wedge *\beta,
\]
wobei $*: \Omega^k \to \Omega^{m-k}$ der durch $g$ bestimmte Hodge-Stern-Operator ist.
Das formale Adjungierte von $d: \Omega^k \to \Omega^{k+1}$ ist
$d^* = (-1)^{mk+m+1} * d\, * : \Omega^k \to \Omega^{k-1}$.

\begin{definition}
Der \emph{Laplace–Beltrami-Operator} (oder Hodge-Laplacian) ist
\[
  \Delta = dd^* + d^*d : \Omega^k(X) \longrightarrow \Omega^k(X).
\]
Eine Form $\alpha$ heißt \emph{harmonisch}, wenn $\Delta\alpha = 0$, oder äquivalent
(für kompaktes $X$): $d\alpha = 0$ und $d^*\alpha = 0$.
\end{definition}

\subsection{Der Hodge-Zerlegungssatz}

\begin{theorem}[Hodge-Satz, bewiesen]
\label{thm:hodge-main}
Sei $X$ eine kompakte orientierte Riemannsche Mannigfaltigkeit. Es gibt eine orthogonale
direkte Summenzerlegung
\[
  \Omega^k(X) = \mathcal{H}^k(X) \oplus \mathrm{im}(d) \oplus \mathrm{im}(d^*),
\]
wobei $\mathcal{H}^k(X) = \ker(\Delta) \cap \Omega^k(X)$ der endlich-dimensionale Raum
der harmonischen $k$-Formen ist. Ferner ist die natürliche Abbildung
\[
  \mathcal{H}^k(X) \xrightarrow{\;\sim\;} H^k_{\mathrm{dR}}(X, \mathbb{R})
\]
ein Isomorphismus.
\end{theorem}

\begin{proof}[Beweisidee]
Die elliptische Regularitätstheorie zeigt, dass $\Delta$ ein selbstadjungierter
elliptischer Operator ist, also kompakte Resolvente hat. Das Fredholm-Prinzip liefert
die orthogonale Zerlegung, und die endliche Dimension von $\mathcal{H}^k$ folgt aus
der Kompaktheit. Der Isomorphismus mit der de-Rham-Kohomologie ist unmittelbar, da
$\ker d = \mathcal{H}^k \oplus \mathrm{im}(d)$.
Siehe \cite[Kap.~0]{griffiths_harris1978} oder \cite[Kap.~4]{warner1983}.
\end{proof}

\subsection{Hodge-Zerlegung für Kähler-Mannigfaltigkeiten}

Auf einer kompakten Kähler-Mannigfaltigkeit liefert die komplexe Struktur eine feinere
Zerlegung.

\begin{theorem}[Hodge-Zerlegung für Kähler-Mannigfaltigkeiten, bewiesen]
\label{thm:hodge-kahler}
Sei $X$ eine kompakte Kähler-Mannigfaltigkeit der komplexen Dimension $n$. Es gibt eine
direkte Summenzerlegung
\[
  H^k(X, \mathbb{C}) = \bigoplus_{p+q=k} H^{p,q}(X),
\]
wobei $H^{p,q}(X)$ den Raum der durch $\bar\partial$-geschlossene $(p,q)$-Formen
darstellbaren Kohomologieklassen (Dolbeault-Kohomologie) bezeichnet. Ferner gilt:
\begin{enumerate}[label=(\roman*)]
\item $\overline{H^{p,q}(X)} = H^{q,p}(X)$, also $h^{p,q} = h^{q,p}$;
\item $h^{p,q} = h^{n-p,n-q}$ nach Serre-Dualität;
\item jedes $H^{p,q}(X)$ wird durch harmonische Formen vom Typ $(p,q)$ repräsentiert.
\end{enumerate}
\end{theorem}

Die ganzen Zahlen $h^{p,q} = \dim_{\mathbb{C}} H^{p,q}(X)$ heißen \emph{Hodge-Zahlen}
und sind topologische Invarianten der Kähler-Mannigfaltigkeit (sie hängen von der
komplexen Struktur ab; biholomerphe Mannigfaltigkeiten haben dieselben Hodge-Zahlen).

\begin{definition}
Das \emph{Hodge-Diamant} einer kompakten Kähler-Mannigfaltigkeit $X$ der komplexen
Dimension $n$ ist die Anordnung $(h^{p,q})_{0 \le p,q \le n}$ in Diamantform:
\[
\begin{array}{ccccc}
  & & h^{n,n} & & \\
  & h^{n,n-1} & & h^{n-1,n} & \\
  & \cdots & & \cdots & \\
  h^{0,0} & & & & \\
\end{array}
\]
\end{definition}

\begin{beispiel}[Hodge-Diamant von $\mathbb{CP}^2$]
Für die komplexe projektive Ebene $X = \mathbb{CP}^2$ (komplexe Dimension $n=2$)
ist der Hodge-Diamant
\[
\begin{array}{ccccc}
 & & 1 & & \\
 & 0 & & 0 & \\
 1 & & 1 & & 1 \\
 & 0 & & 0 & \\
 & & 1 & &
\end{array}
\]
mit $h^{0,0}=h^{1,1}=h^{2,2}=1$ und allen anderen $h^{p,q}=0$.
\end{beispiel}

%% ============================================================
\section{de-Rham- versus Dolbeault-Kohomologie}
\label{sec:kohomologie}
%% ============================================================

\subsection{de-Rham-Kohomologie}

Die de-Rham-Kohomologiegruppen
\[
  H^k_{\mathrm{dR}}(X, \mathbb{R}) = \frac{\ker(d: \Omega^k \to \Omega^{k+1})}{\mathrm{im}(d: \Omega^{k-1} \to \Omega^k)}
\]
sind rein topologische Invarianten: Durch den de-Rham-Satz (bewiesen) sind sie
isomorph zur singulären Kohomologie $H^k(X, \mathbb{R})$.

\subsection{Dolbeault-Kohomologie}

Auf einer komplexen Mannigfaltigkeit zerfällt das äußere Differential als $d = \partial + \bar\partial$,
wobei $\partial: \mathcal{A}^{p,q} \to \mathcal{A}^{p+1,q}$ und
$\bar\partial: \mathcal{A}^{p,q} \to \mathcal{A}^{p,q+1}$. Die \emph{Dolbeault-Kohomologiegruppen}
sind
\[
  H^{p,q}(X) = H^{p,q}_{\bar\partial}(X) = \frac{\ker(\bar\partial: \mathcal{A}^{p,q} \to \mathcal{A}^{p,q+1})}{\mathrm{im}(\bar\partial: \mathcal{A}^{p,q-1} \to \mathcal{A}^{p,q})}.
\]
Durch den Dolbeault-Satz (bewiesen) gilt $H^{p,q}(X) \cong H^q(X, \Omega^p_X)$,
wobei $\Omega^p_X$ die Garbe der holomorphen $p$-Formen bezeichnet.

\subsection{Die Hodge-zu-de-Rham-Spektralsequenz}

Es gibt eine Spektralsequenz (die Hodge-zu-de-Rham- oder Frölicher-Spektralsequenz)
mit $E_1^{p,q} = H^q(X, \Omega^p_X)$, die gegen $H^{p+q}(X, \mathbb{C})$ konvergiert.
Für Kähler-Mannigfaltigkeiten degeneriert diese Spektralsequenz bei $E_1$ (bewiesen),
was äquivalent zum Hodge-Zerlegungssatz ist. Für allgemeine komplexe Mannigfaltigkeiten
muss das nicht gelten.

\subsection{Die Hodge-Filtrierung}

Die Hodge-Zerlegung definiert eine absteigende Filtrierung
\[
  F^p H^k(X, \mathbb{C}) = \bigoplus_{a \ge p} H^{a,k-a}(X),
\]
die \emph{Hodge-Filtrierung}. Sie ist funktoriell und zentral in der abstrakten
Theorie der Hodge-Strukturen.

\begin{definition}
Eine \emph{reine Hodge-Struktur vom Gewicht $k$} ist ein endlich-dimensionaler reeller
Vektorraum $H_\mathbb{R}$ zusammen mit einer Bgraduierung
$H_\mathbb{C} = \bigoplus_{p+q=k} H^{p,q}$, die $\overline{H^{p,q}} = H^{q,p}$ erfüllt.
\end{definition}

%% ============================================================
\section{Algebraische Zyklen}
\label{sec:zyklen}
%% ============================================================

\subsection{Algebraische Teilvarietäten und ihre Klassen}

Sei $X$ eine glatte projektive komplexe algebraische Varietät der komplexen Dimension $n$.

\begin{definition}
Ein \emph{algebraischer Zyklus der Kodimension $k$} auf $X$ ist eine formale
ganzzahlige Linearkombination
\[
  Z = \sum_{i} n_i Z_i, \quad n_i \in \mathbb{Z},
\]
wobei jedes $Z_i \subset X$ eine irreduzible abgeschlossene Teilvarietät der Kodimension $k$ ist.
\end{definition}

\begin{definition}
Die \emph{Chow-Gruppe} $\mathrm{CH}^k(X)$ ist die Gruppe der algebraischen Zyklen
der Kodimension $k$ modulo rationaler Äquivalenz.
\end{definition}

\subsection{Die Zyklenabbildung und Fundamentalklassen}

Jede irreduzible Teilvarietät $Z \subset X$ der komplexen Kodimension $k$ (also
reellen Kodimension $2k$) trägt eine Fundamentalklasse $[Z] \in H_{2n-2k}(X, \mathbb{Z})$.
Durch Poincaré-Dualität entspricht dies einer Kohomologieklasse
\[
  [Z]^* \in H^{2k}(X, \mathbb{Z}).
\]

\begin{theorem}[Typ von Zyklenklassen, bewiesen]
\label{thm:zyklus-typ}
Die Kohomologieklasse $[Z]^* \in H^{2k}(X, \mathbb{Z}) \subset H^{2k}(X, \mathbb{C})$
liegt unter der Hodge-Zerlegung in $H^{k,k}(X)$.
\end{theorem}

\begin{proof}[Beweisidee]
Dies folgt daraus, dass $Z$ lokal durch holomorphe Gleichungen definiert ist: Der
Poincaré-duale Strom $[Z]$ ist vom Typ $(k,k)$, weil er auf einer komplexen
Untermannigfaltigkeit getragen wird (glatter Fall), oder durch Auflösung von
Singularitäten (allgemeiner Fall).
\end{proof}

\begin{definition}
Das Bild der \emph{Zyklenklassenabbildung}
\[
  \mathrm{kl}: \mathrm{CH}^k(X) \longrightarrow H^{2k}(X, \mathbb{Z})
\]
besteht aus \emph{algebraischen Kohomologieklassen}. Eine Klasse
$\alpha \in H^{2k}(X,\mathbb{Z})$, die in $H^{k,k}(X)$ liegt, heißt
\emph{ganzzahlige Hodge-Klasse}, und eine Klasse in
$H^{2k}(X, \mathbb{Q}) \cap H^{k,k}(X)$ heißt \emph{rationale Hodge-Klasse}.
\end{definition}

%% ============================================================
\section{Die Hodge-Vermutung}
\label{sec:vermutung}
%% ============================================================

\subsection{Formale Aussage}

Wir sind nun in der Lage, die Hodge-Vermutung präzise zu formulieren.

\begin{vermutung}[Hodge-Vermutung, offen]
\label{verm:hodge}
Sei $X$ eine glatte projektive komplexe algebraische Varietät. Für jedes $k \ge 0$
induziert die Zyklenklassenabbildung eine Surjektion
\[
  \mathrm{kl}_\mathbb{Q}: \mathrm{CH}^k(X) \otimes_\mathbb{Z} \mathbb{Q}
  \longrightarrow
  H^{2k}(X, \mathbb{Q}) \cap H^{k,k}(X).
\]
Äquivalent: Jede rationale Kohomologieklasse vom Typ $(k,k)$ ist eine rationale
Linearkombination von Klassen algebraischer Teilvarietäten.
\end{vermutung}

\begin{bemerkung}
\label{bem:millennium}
Die Hodge-Vermutung ist eines der sieben Millennium-Preisprobleme des Clay Mathematics
Institute \cite{clay2000}. Sie ist seit Hodges ursprünglicher Formulierung im Jahr
1950 offen \cite{hodge1950}. Stand 2026 ist sie für alle $k \ge 2$ auf Varietäten
der Dimension $n \ge 3$ unbewiesen.
\end{bemerkung}

\subsection{Äquivalente Umformulierungen}

\begin{proposition}[Äquivalente Form, bewiesen]
Die folgenden Aussagen sind für eine glatte projektive Varietät $X$ äquivalent:
\begin{enumerate}[label=(\roman*)]
\item Jede rationale $(k,k)$-Hodge-Klasse auf $X$ ist algebraisch.
\item Die Zyklenabbildung $\mathrm{CH}^k(X)_\mathbb{Q} \to H^{k,k}(X) \cap H^{2k}(X,\mathbb{Q})$
  ist surjektiv.
\item Das natürliche Bild von $\mathrm{Hdg}^k(X)$ in
  $H^{2k}(X, \mathbb{Q}) / N^{k+1} H^{2k}(X,\mathbb{Q})$ wird von algebraischen
  Klassen erzeugt.
\end{enumerate}
\end{proposition}

\subsection{Warum rationale Koeffizienten?}

Hodges ursprüngliche Formulierung von 1950 verwendete \emph{ganzzahlige} Koeffizienten.
Wie wir in Abschnitt~\ref{sec:gegenbeispiele} diskutieren, ist die ganzzahlige Version
\emph{falsch} — es gibt ganzzahlige Hodge-Klassen, die keine ganzzahligen Linearkombinationen
von Zyklenklassen sind. Die korrekte Version erfordert rationale Koeffizienten.

\subsection{Funktorialität}

\begin{proposition}[Funktorialität, bewiesen]
Wenn $f: X \to Y$ ein Morphismus glatter projektiver Varietäten ist, dann bildet
$f^*$ Hodge-Klassen auf $Y$ auf Hodge-Klassen auf $X$ ab, und $f_*$ (wo definiert)
bildet algebraische Klassen auf algebraische Klassen ab.
\end{proposition}

%% ============================================================
\section{Bekannte Fälle}
\label{sec:bekannt}
%% ============================================================

Trotz ihrer scheinbaren Einfachheit ist die Hodge-Vermutung in bemerkenswert wenigen
Fällen bekannt.

\subsection{Triviale Fälle}

\begin{itemize}
\item $k = 0$: $H^0(X,\mathbb{Q}) \cap H^{0,0}(X) = \mathbb{Q}$, erzeugt durch
  die Fundamentalklasse $[X]$. Trivialerweise wahr.
\item $k = n$: Durch Serre-Dualität äquivalent zum Fall $k = 0$ auf der dualen
  Varietät, also trivialerweise wahr.
\item $n = 1$ (Kurven): $X$ ist eine Riemann-Fläche; nur $H^{0,0}$ und $H^{1,1}$
  kommen vor, beide trivial.
\end{itemize}

\subsection{Der Lefschetz-Fall $k=1$}

Der erste nicht-triviale Fall ist $k = 1$, der durch den Lefschetz-$(1,1)$-Satz
behandelt wird (Abschnitt~\ref{sec:lefschetz}, bewiesen von Lefschetz 1924 und
vervollständigt durch Kodaira).

\subsection{Abelsche Varietäten}

Für abelsche Varietäten $A$ der Dimension $n$ gilt für die Hodge-Vermutung:
\begin{itemize}
\item Bewiesen für $k = 1$ (Lefschetz-Satz).
\item Bewiesen für $k = n-1$ (durch Dualität mit $k = 1$).
\item Bekannt für gewisse spezielle abelsche Varietäten: Weil-4-Faltigkeiten
  (partiell), Potenzen elliptischer Kurven (Tate, Murasaki), einfache abelsche
  Varietäten primer Dimension (für $k=1$ via Albert-Klassifikation).
\item Offen für $k = 2$ auf allgemeinen abelschen 4-Faltigkeiten.
\end{itemize}

\subsection{Unirulierte und rational zusammenhängende Varietäten}

Für rational zusammenhängende glatte projektive Varietäten $X$ ist die
Hodge-Vermutung in Kodimension 1 bekannt (Lefschetz), bleibt aber in höherer
Kodimension offen. Die Hodge-Zahlen rational zusammenhängender Varietäten
erfüllen $h^{p,0} = 0$ für $p > 0$ (bewiesen).

\subsection{Die Tate-Vermutung}

Über endlichen Körpern gibt es ein Analogon: Die \emph{Tate-Vermutung} besagt,
dass Zyklenklassen die (kristalline oder $\ell$-adische) Kohomologie in Grad $2k$
erzeugen. Tate-Vermutung und Hodge-Vermutung sind eng verwandt, ohne dass eine
die andere impliziert.

%% ============================================================
\section{Der Lefschetz-\texorpdfstring{$(1,1)$}{(1,1)}-Satz}
\label{sec:lefschetz}
%% ============================================================

Der Lefschetz-$(1,1)$-Satz ist das Prunkstück der bewiesenen Fälle der Hodge-Vermutung.
Er charakterisiert algebraische Kohomologieklassen in Grad 2 vollständig.

\subsection{Aussage und Beweis}

\begin{theorem}[Lefschetz-$(1,1)$-Satz, bewiesen]
\label{thm:lefschetz11}
Sei $X$ eine kompakte Kähler-Mannigfaltigkeit. Eine Klasse $\alpha \in H^2(X, \mathbb{Z})$
liegt in $H^{1,1}(X)$ genau dann, wenn $\alpha$ die erste Chern-Klasse $c_1(L)$
eines holomorphen Geradenbündels $L$ auf $X$ ist.
\end{theorem}

\begin{proof}
Betrachte die exponentielle Garbensequenz auf $X$:
\[
  0 \longrightarrow \mathbb{Z} \xrightarrow{\;\cdot 2\pi i\;}
  \mathcal{O}_X \xrightarrow{\;\exp\;} \mathcal{O}_X^* \longrightarrow 0.
\]
Die zugehörige lange exakte Kohomologiesequenz enthält
\[
  H^1(X, \mathcal{O}_X^*) \xrightarrow{\;c_1\;} H^2(X, \mathbb{Z})
  \xrightarrow{\;\delta\;} H^2(X, \mathcal{O}_X).
\]
Dabei klassifiziert $H^1(X, \mathcal{O}_X^*)$ holomorphe Geradenbündel (die Picard-Gruppe
$\mathrm{Pic}(X)$). Eine Klasse $\alpha \in H^2(X,\mathbb{Z})$ liegt im Bild von $c_1$
genau dann, wenn ihr Bild in $H^2(X, \mathcal{O}_X) = H^{0,2}(X)$ verschwindet.
Nach der Hodge-Zerlegung gilt $H^2(X,\mathbb{C}) = H^{2,0}(X) \oplus H^{1,1}(X) \oplus H^{0,2}(X)$,
und $\alpha \in H^2(X,\mathbb{Z}) \subset H^2(X,\mathbb{C})$ ist reell, also sind seine
$(2,0)$- und $(0,2)$-Komponenten komplex konjugiert. Mithin liegt $\alpha$ in $H^{1,1}(X)$
genau dann, wenn seine $H^{0,2}$-Komponente verschwindet, d.h. genau dann, wenn
$\alpha = c_1(L)$ für ein $L \in \mathrm{Pic}(X)$.
\end{proof}

\begin{korollar}[Hodge-Vermutung in Kodimension 1, bewiesen]
\label{kor:hv1}
Die Hodge-Vermutung gilt für $k = 1$ auf jeder glatten projektiven komplexen Varietät:
Jede rationale $(1,1)$-Hodge-Klasse ist eine rationale Linearkombination von Divisorklassen.
\end{korollar}

\begin{proof}
Nach dem Lefschetz-Satz ist jede ganzzahlige $(1,1)$-Klasse gleich $c_1(L)$ für ein
Geradenbündel $L$. Auf einer projektiven Varietät ist $L^{\otimes m}$ für großes $m$
sehr ampel, entspricht also einem Hyperebenschnitt und damit einer algebraischen Hyperfläche.
Rationale Linearkombinationen liefern das Ergebnis.
\end{proof}

\subsection{Deligne's Beitrag}

Aufbauend auf Grothendieck's Rahmen und seiner eigenen Arbeit zu den Weil-Vermutungen
bewies Deligne in \emph{Hodge III} \cite{deligne1974}, dass die Hodge-de-Rham-Spektralsequenz
für glatte projektive Varietäten auch in gemischten Charakteristiken bei $E_1$ degeneriert.
Dies beweist zwar nicht die Hodge-Vermutung direkt, etabliert aber die Algebraizität von
Hodge-Klassen in gewissen funktoriellen Situationen.

\subsection{Der Hodge-Index-Satz}

\begin{theorem}[Hodge-Index-Satz, bewiesen]
\label{thm:hodge-index}
Sei $X$ eine glatte projektive Fläche, und $H$ eine ample Divisorklasse. Wenn
$D \in \mathrm{NS}(X)$ die Bedingung $D \cdot H = 0$ erfüllt, dann gilt $D^2 \le 0$,
mit Gleichheit genau dann, wenn $D$ numerisch trivial ist.
\end{theorem}

Dieser Satz, bewiesen mit der Hodge-Zerlegung, ist grundlegend für das Studium der
Néron–Severi-Gruppe und von Flächen.

%% ============================================================
\section{Gegenbeispiele und Nicht-Kähler-Mannigfaltigkeiten}
\label{sec:gegenbeispiele}
%% ============================================================

\subsection{Die ganzzahlige Hodge-Vermutung ist falsch}

\begin{theorem}[Atiyah–Hirzebruch, bewiesen]
\label{thm:ah-gegenbeispiel}
Die ganzzahlige Version der Hodge-Vermutung ist falsch: Es gibt glatte projektive
komplexe Varietäten $X$ und ganzzahlige Hodge-Klassen
$\alpha \in H^{2k}(X, \mathbb{Z}) \cap H^{k,k}(X)$, die keine ganzzahligen
Linearkombinationen von Klassen algebraischer Zyklen sind.
\end{theorem}

\begin{proof}[Historischer Kontext]
Atiyah und Hirzebruch \cite{atiyah_hirzebruch1961} konstruierten das erste Gegenbeispiel
unter Verwendung von Torsionsklassen in der Kohomologie gewisser Varietäten. Ihre Methode
nutzt die Atiyah–Hirzebruch-Spektralsequenz und zeigt, dass Torsions-Hodge-Klassen nicht
von algebraischen Zyklen stammen müssen.
\end{proof}

\subsection{Voisins stärkeres Gegenbeispiel}

\begin{theorem}[Voisin 2002, bewiesen]
\label{thm:voisin}
Es gibt glatte projektive komplexe Varietäten $X$ und rationale Hodge-Klassen
$\alpha \in H^{2k}(X, \mathbb{Q}) \cap H^{k,k}(X)$, so dass $\alpha$ auch nach
Übergang zu einem birationalen Modell nicht die Klasse eines algebraischen Zyklus ist.
\end{theorem}

\begin{bemerkung}
Voisins Ergebnis \cite{voisin2002} ist stärker als Atiyah–Hirzebruch: Es zeigt, dass
sogar bis auf birationale Äquivalenz die ganzzahlige Hodge-Vermutung versagt.
Allerdings liefert es kein Gegenbeispiel zur \emph{rationalen} Hodge-Vermutung
(Vermutung~\ref{verm:hodge}), die offen bleibt.
\end{bemerkung}

\subsection{Nicht-Kähler-Mannigfaltigkeiten}

Die Hodge-Vermutung ist für projektive (also Kähler-) Varietäten formuliert.
Für allgemeine komplexe Mannigfaltigkeiten kann die Hodge-Zerlegung selbst versagen:

\begin{theorem}[bewiesen]
Es gibt kompakte komplexe Mannigfaltigkeiten, die nicht Kähler sind (z.B. Hopf-Flächen
$(\mathbb{C}^2 \setminus \{0\})/\langle z \mapsto 2z \rangle$), für die der
Hodge-Zerlegungssatz nicht gilt und für die jedes Analogon der Hodge-Vermutung
leicht scheitert.
\end{theorem}

\begin{bemerkung}
Die Kähler-(bzw. projektiv-)Voraussetzung in der Hodge-Vermutung ist also wesentlich
und kann nicht fallen gelassen werden.
\end{bemerkung}

\subsection{Hodge-Vermutung für singuläre Varietäten}

Die Hodge-Vermutung in der angegebenen Form erfordert, dass $X$ glatt ist. Für
singuläre Varietäten verwendet man Schnitt-Kohomologie (Goresky–MacPherson) oder
gemischte Hodge-Strukturen (Deligne), aber die entsprechenden Vermutungen sind
subtiler und in der Literatur nicht vollständig ausformuliert.

%% ============================================================
\section{Motive und Standard-Vermutungen}
\label{sec:motive}
%% ============================================================

\subsection{Grothendieck's Theorie der Motive}

In den 1960er Jahren schlug Grothendieck eine universelle Kohomologietheorie vor —
die Theorie der \emph{Motive} — die de-Rham-, Betti-, $\ell$-adische und kristalline
Kohomologien als spezielle Realisierungen umfassen sollte \cite{grothendieck1969}.

\begin{definition}
Die \emph{Kategorie der reinen Motive} $\mathbf{Mot}(k)_\mathbb{Q}$ über einem
Körper $k$ (mit rationalen Koeffizienten) hat:
\begin{itemize}
\item Objekte: Tripel $(X, p, n)$, wobei $X$ eine glatte projektive Varietät,
  $p \in \mathrm{CH}^{\dim X}(X \times X)_\mathbb{Q}$ ein idempotenter Korrespondenz
  und $n \in \mathbb{Z}$ eine Tate-Verdrehung ist.
\item Morphismen: Korrespondenzen modulo rationaler Äquivalenz.
\end{itemize}
\end{definition}

Die Hodge-Vermutung ist äquivalent zur Behauptung, dass jede Hodge-Klasse
\emph{motiviert} ist, d.h. im Bild der motivischen Zyklenklassenabbildung liegt.

\subsection{Die Standard-Vermutungen}

Grothendieck formulierte vier \emph{Standard-Vermutungen} (A, B, C, D), die, wenn
bewiesen, die Kategorie der Motive als halbeinfache abelsche Kategorie etablieren
und die Weil-Vermutungen (von Deligne ohne sie bewiesen) implizieren würden:

\begin{vermutung}[Lefschetz-Standard-Vermutung B, offen]
\label{verm:std-b}
Für eine glatte projektive Varietät $X$ über einem algebraisch abgeschlossenen Körper
und eine ample Klasse $L$ sind die Lefschetz-Operatoren $\Lambda^k$ (Inverse zu $L^k$
auf primitiver Kohomologie) algebraisch, d.h. durch algebraische Korrespondenzen gegeben.
\end{vermutung}

\begin{vermutung}[Künneth-Standard-Vermutung C, offen]
\label{verm:std-c}
Die Künneth-Projektoren $\pi_k: H^*(X) \to H^k(X)$ sind algebraisch.
\end{vermutung}

\begin{vermutung}[Hodge-Standard-Vermutung D, offen]
\label{verm:std-d}
Numerische und homologische Äquivalenz algebraischer Zyklen fallen zusammen.
\end{vermutung}

\begin{bemerkung}
Die Standard-Vermutungen sind über allen Körpern, einschließlich $\mathbb{C}$, offen.
Vermutung~\ref{verm:std-b} über $\mathbb{C}$ würde aus der Hodge-Vermutung folgen,
ist aber anderweitig unbekannt. Vermutung~\ref{verm:std-d} ist für abelsche Varietäten
bewiesen (Lieberman 1968).
\end{bemerkung}

\subsection{Verbindung zu L-Funktionen}

Motive tragen natürliche L-Funktionen. Für eine glatte projektive Varietät $X/\mathbb{Q}$
ist die \emph{motivische L-Funktion}
\[
  L(H^k(X), s) = \prod_p \det\bigl(1 - p^{-s} \mathrm{Frob}_p \,|\, H^k(\bar X)^{I_p}\bigr)^{-1}
\]
für eine Funktionalgleichung und analytische Fortsetzung erwartet. Die
Tate-Vermutung sagt voraus, dass die Ordnung des Pols bei $s = k/2 + 1$ den Rang von
$\mathrm{CH}^{k/2}(X)$ gleicht. Die BSD-Vermutung ist der Fall $k = 1$ für elliptische Kurven.

\subsection{Grothendieck's Vision}

Grothendieck glaubte, die Hodge-Vermutung würde aus den Standard-Vermutungen via
der Existenz einer motivischen Galois-Gruppe folgen:

\begin{vermutung}[Motivische Galois-Gruppe, offen]
\label{verm:motivische-galois}
Die Tannaka-Kategorie der reinen Motive über $\mathbb{C}$ ist äquivalent zur
Kategorie der Darstellungen einer pro-reduktiven algebraischen Gruppe $G_{\mathrm{mot}}$,
der \emph{motivischen Galois-Gruppe}. Die Hodge-Vermutung würde aus der Algebraizität
der Hodge-Gruppe $\mathrm{MT}(H) \subset G_{\mathrm{mot}}$ folgen.
\end{vermutung}

%% ============================================================
\section*{Schluss}
\label{sec:schluss}
%% ============================================================

Die Hodge-Vermutung bleibt eines der tiefsten ungelösten Probleme der Mathematik.
Wir haben gesehen:

\begin{enumerate}[label=(\roman*)]
\item Der \textbf{Hodge-Zerlegungssatz} (bewiesen) und das $\partial\bar\partial$-Lemma
  (bewiesen) sind Eckpfeiler der Kähler-Geometrie.
\item Klassen algebraischer Zyklen sind stets vom Hodge-Typ $(k,k)$ (bewiesen);
  die Hodge-Vermutung fragt nach der Umkehrung.
\item Der \textbf{Lefschetz-$(1,1)$-Satz} (bewiesen) klärt den Fall $k=1$:
  Ganzzahlige $(1,1)$-Klassen sind genau Chern-Klassen von Geradenbündeln.
\item Die \textbf{ganzzahlige Hodge-Vermutung ist falsch} (Atiyah–Hirzebruch,
  bewiesen; Voisin 2002, bewiesen), was die rationale Version als die korrekte
  Formulierung auszeichnet.
\item Die rationale Hodge-Vermutung ist für alle $k \ge 2$ auf Varietäten der
  Dimension $\ge 3$ offen.
\item Grothendieck's Theorie der Motive liefert einen konjekcturellen Rahmen
  (Standard-Vermutungen, alle offen), der die Hodge-Vermutung implizieren würde.
\end{enumerate}

Jeder Fortschritt bei der Hodge-Vermutung würde einen Durchbruch darstellen, der
die drei großen Säulen der Geometrie verbindet: Topologie, komplexe Analysis und
algebraische Geometrie.

%% ============================================================
\section{Literatur}
\label{sec:literatur}
%% ============================================================

\begin{thebibliography}{99}

\bibitem{hodge1950}
W.~V.~D. Hodge,
\textit{The topological invariants of algebraic varieties},
Proc. Int. Congr. Math. (Cambridge, MA, 1950), Bd.~1, 182--192.

\bibitem{lefschetz1924}
S.~Lefschetz,
\textit{L'analysis situs et la g\'eom\'etrie alg\'ebrique},
Gauthier-Villars, Paris, 1924.

\bibitem{griffiths_harris1978}
P.~Griffiths und J.~Harris,
\textit{Principles of Algebraic Geometry},
Wiley-Interscience, New York, 1978.

\bibitem{deligne1974}
P.~Deligne,
\textit{Th\'eorie de Hodge: III},
Publ. Math. IHES \textbf{44} (1974), 5--77.

\bibitem{voisin2002}
C.~Voisin,
\textit{A counterexample to the Hodge conjecture extended to K\"ahler varieties},
Int. Math. Res. Not. \textbf{2002}, Nr.~20, 1057--1075.

\bibitem{grothendieck1969}
A.~Grothendieck,
\textit{Standard conjectures on algebraic cycles},
Algebraic Geometry (Bombay Colloquium 1968), Oxford Univ. Press, 1969, 193--199.

\bibitem{atiyah_hirzebruch1961}
M.~F. Atiyah und F.~Hirzebruch,
\textit{Analytic cycles on complex manifolds},
Topology \textbf{1} (1962), 25--45.

\bibitem{weil1949}
A.~Weil,
\textit{Numbers of solutions of equations in finite fields},
Bull. Amer. Math. Soc. \textbf{55} (1949), 497--508.

\bibitem{voisin2007}
C.~Voisin,
\textit{Hodge Theory and Complex Algebraic Geometry}, Bde.~I--II,
Cambridge Univ. Press, 2002--2003.

\bibitem{clay2000}
A.~Jaffe und A.~Wiles (Hrsg.),
\textit{The Millennium Prize Problems},
Clay Mathematics Institute, Cambridge, MA, 2006.

\bibitem{warner1983}
F.~W. Warner,
\textit{Foundations of Differentiable Manifolds and Lie Groups},
Springer, New York, 1983.

\end{thebibliography}

\end{document}
